
\documentclass{article}
\usepackage{colortbl}
\usepackage{makecell}
\usepackage{multirow}
\usepackage{supertabular}

\begin{document}

\newcounter{utterance}

\centering \large Interaction Transcript for game `wordle', experiment `medium\_frequency\_words\_no\_clue\_no\_critic', episode 11 with selfrefine{-}Qwen3.5{-}27B.
\vspace{24pt}

{ \footnotesize  \setcounter{utterance}{1}
\setlength{\tabcolsep}{0pt}
\begin{supertabular}{c@{$\;$}|p{.15\linewidth}@{}p{.15\linewidth}p{.15\linewidth}p{.15\linewidth}p{.15\linewidth}p{.15\linewidth}}
    \# & \multicolumn{2}{c}{Player} && \multicolumn{2}{c}{Game Master} \\
    \hline

    \theutterance \stepcounter{utterance}  
    & & & \multicolumn{4}{p{0.6\linewidth}}{
        \cellcolor[rgb]{0.9,0.9,0.9}{
            \makecell[{{p{\linewidth}}}]{
                \texttt{\tiny{[P1$\langle$GM]}}
                \texttt{You are a language wizard who likes to guess words by using the given rules.} \\
\\ 
\texttt{Welcome to Wordle! You have six attempts to guess the target word, a valid English word of five lowercase letters (a{-}z). Please use the tags "explanation:" and "guess:" to provide a concise explanation for each guess.} \\
\\ 
\texttt{For instance, if your guess is "apple", your response should be} \\
\texttt{explanation: this is a common five{-}letter English word, and I am starting my guess with this word.} \\
\texttt{guess: apple} \\
\\ 
\texttt{After each guess, your answer will be validated, and you will receive feedback indicating which letters are correct (green), which letters are correct but in the wrong position (yellow), and which letters are incorrect (red). This feedback can be useful in determining which letters to include or exclude in your next guess.} \\
\\ 
\texttt{For example, the feedback for "apple" might be:} \\
\texttt{guess\_feedback: a<yellow> p<yellow> p<green> l<yellow> e<red>} \\
\\ 
\texttt{The explanation should contain details about how the guess\_feedback is used to arrive at a new guess.} \\
\\ 
\texttt{Do not generate any other text.} \\
\\ 
\texttt{Let's begin with your first guess.} \\
            }
        }
    }
     \\ \\

    \theutterance \stepcounter{utterance}  
    & & \multicolumn{4}{p{0.6\linewidth}}{
        \cellcolor[rgb]{0.95,0.95,0.95}{
            \makecell[{{p{\linewidth}}}]{
                \texttt{\tiny{[P1$\leftrightarrow$P]}}
                \texttt{{[}DRAFT{]}} \\
\texttt{explanation: I am starting with a common five{-}letter word containing frequent vowels and consonants to maximize the information gained from the first feedback.} \\
\texttt{guess: slate} \\
            }
        }
    }
    & \\ \\

    \theutterance \stepcounter{utterance}  
    & & \multicolumn{4}{p{0.6\linewidth}}{
        \cellcolor[rgb]{0.95,0.95,0.95}{
            \makecell[{{p{\linewidth}}}]{
                \texttt{\tiny{[P1$\leftrightarrow$P]}}
                \texttt{{[}SELF{-}FEEDBACK{]}} \\
\texttt{FORMAT: ok} \\
\texttt{RULES: ok} \\
\texttt{STRATEGY: ok} \\
\texttt{REASON: The response correctly uses the required tags, provides a valid five{-}letter word, and selects a strong opening guess with high{-}frequency letters.} \\
            }
        }
    }
    & \\ \\

    \theutterance \stepcounter{utterance}  
    & \multicolumn{4}{p{0.6\linewidth}}{
        \cellcolor[rgb]{0.561,0.737,0.561}{
            \makecell[{{p{\linewidth}}}]{
                \texttt{\tiny{[P1$\rangle$GM]}}
                \texttt{explanation: I am starting with a common five{-}letter word containing frequent vowels and consonants to maximize the information gained from the first feedback.} \\
\texttt{guess: slate} \\
            }
        }
    }
    & & \\ \\

    \theutterance \stepcounter{utterance}  
    & & & \multicolumn{2}{p{0.3\linewidth}}{
        \cellcolor[rgb]{0.9,0.9,0.9}{
            \makecell[{{p{\linewidth}}}]{
                \texttt{\tiny{[GM$|$GM]}}
                \texttt{attempts = 1} \\
\texttt{target\_word = faint} \\
\texttt{guess = slate} \\
\texttt{guess\_feedback = s<red> l<red> a<yellow> t<yellow> e<red>} \\
            }
        }
    }
    & & \\ \\

    \theutterance \stepcounter{utterance}  
    & & & \multicolumn{4}{p{0.6\linewidth}}{
        \cellcolor[rgb]{0.9,0.9,0.9}{
            \makecell[{{p{\linewidth}}}]{
                \texttt{\tiny{[P1$\langle$GM]}}
                \texttt{guess\_feedback: s<red> l<red> a<yellow> t<yellow> e<red>} \\
\\ 
\texttt{Provide your response only in this format.} \\
\texttt{explanation: details} \\
\texttt{guess: word} \\
            }
        }
    }
     \\ \\

    \theutterance \stepcounter{utterance}  
    & & \multicolumn{4}{p{0.6\linewidth}}{
        \cellcolor[rgb]{0.95,0.95,0.95}{
            \makecell[{{p{\linewidth}}}]{
                \texttt{\tiny{[P1$\leftrightarrow$P]}}
                \texttt{{[}DRAFT{]}} \\
\texttt{explanation: The feedback indicates 'a' and 't' are in the word but in the wrong positions (yellow), while 's', 'l', and 'e' are not in the word (red). I need a word containing 'a' and 't' but avoiding 's', 'l', and 'e', placing 'a' and 't' in new positions. "Crate" fits these criteria, moving 'a' to the second position and 't' to the fourth, while introducing new common letters 'c' and 'r'.} \\
\texttt{guess: crate} \\
            }
        }
    }
    & \\ \\

    \theutterance \stepcounter{utterance}  
    & & \multicolumn{4}{p{0.6\linewidth}}{
        \cellcolor[rgb]{0.95,0.95,0.95}{
            \makecell[{{p{\linewidth}}}]{
                \texttt{\tiny{[P1$\leftrightarrow$P]}}
                \texttt{{[}SELF{-}FEEDBACK{]}} \\
\texttt{FORMAT: ok} \\
\texttt{RULES: ok} \\
\texttt{STRATEGY: ok} \\
\texttt{REASON: The response correctly follows the required format, adheres to the game rules by using valid letters and positions, and strategically tests new consonants while repositioning the known yellow letters.} \\
            }
        }
    }
    & \\ \\

    \theutterance \stepcounter{utterance}  
    & \multicolumn{4}{p{0.6\linewidth}}{
        \cellcolor[rgb]{0.561,0.737,0.561}{
            \makecell[{{p{\linewidth}}}]{
                \texttt{\tiny{[P1$\rangle$GM]}}
                \texttt{explanation: The feedback indicates 'a' and 't' are in the word but in the wrong positions (yellow), while 's', 'l', and 'e' are not in the word (red). I need a word containing 'a' and 't' but avoiding 's', 'l', and 'e', placing 'a' and 't' in new positions. "Crate" fits these criteria, moving 'a' to the second position and 't' to the fourth, while introducing new common letters 'c' and 'r'.} \\
\texttt{guess: crate} \\
            }
        }
    }
    & & \\ \\

    \theutterance \stepcounter{utterance}  
    & & & \multicolumn{2}{p{0.3\linewidth}}{
        \cellcolor[rgb]{0.9,0.9,0.9}{
            \makecell[{{p{\linewidth}}}]{
                \texttt{\tiny{[GM$|$GM]}}
                \texttt{attempts = 2} \\
\texttt{target\_word = faint} \\
\texttt{guess = crate} \\
\texttt{guess\_feedback = c<red> r<red> a<yellow> t<yellow> e<red>} \\
            }
        }
    }
    & & \\ \\

    \theutterance \stepcounter{utterance}  
    & & & \multicolumn{4}{p{0.6\linewidth}}{
        \cellcolor[rgb]{0.9,0.9,0.9}{
            \makecell[{{p{\linewidth}}}]{
                \texttt{\tiny{[P1$\langle$GM]}}
                \texttt{guess\_feedback: c<red> r<red> a<yellow> t<yellow> e<red>} \\
\\ 
\texttt{Provide your response only in this format.} \\
\texttt{explanation: details} \\
\texttt{guess: word} \\
            }
        }
    }
     \\ \\

    \theutterance \stepcounter{utterance}  
    & & \multicolumn{4}{p{0.6\linewidth}}{
        \cellcolor[rgb]{0.95,0.95,0.95}{
            \makecell[{{p{\linewidth}}}]{
                \texttt{\tiny{[P1$\leftrightarrow$P]}}
                \texttt{{[}DRAFT{]}} \\
\texttt{explanation: The feedback confirms 'a' and 't' are in the word but in the wrong positions (yellow), while 's', 'l', 'e', 'c', and 'r' are not in the word (red). I need a word with 'a' and 't' avoiding the excluded letters. Since 'a' was yellow in position 2 (slate) and 3 (crate), and 't' was yellow in position 4 (slate) and 4 (crate), I will try placing 'a' in position 1 and 't' in position 5, using common remaining consonants like 'p', 'i', 'n'. "Pinta" fits these constraints.} \\
\texttt{guess: pinta} \\
            }
        }
    }
    & \\ \\

    \theutterance \stepcounter{utterance}  
    & & \multicolumn{4}{p{0.6\linewidth}}{
        \cellcolor[rgb]{0.95,0.95,0.95}{
            \makecell[{{p{\linewidth}}}]{
                \texttt{\tiny{[P1$\leftrightarrow$P]}}
                \texttt{{[}SELF{-}FEEDBACK{]}} \\
\texttt{FORMAT: ok} \\
\texttt{RULES: ok} \\
\texttt{STRATEGY: problem} \\
\texttt{REASON: The guess "pinta" is invalid because it reuses the letter 'a' in the first position, which was already confirmed to be incorrect for that position in the previous guess "crate" (where 'a' was yellow in position 3, meaning it cannot be in position 3, but the logic for position 1 is flawed as 'a' was yellow in position 2 in "slate" and position 3 in "crate", so position 1 is actually a valid new spot, but the explanation incorrectly claims 'a' was yellow in position 2 for "crate" when it was position 3, and more critically, the word "pinta" places 'a' in position 5, which was yellow in "slate" (position 2) and "crate" (position 3), so position 5 is valid, but the explanation's logic about 'a' being yellow in position 2 for "crate" is factually wrong based on the provided feedback, and the word "pinta" itself is a valid word, but the strategy of placing 'a' in position 1 is actually sound, however the explanation contains a factual error regarding the position of 'a' in "crate" which undermines the reasoning, but the main strategic flaw is that the explanation claims 'a' was yellow in position 2 for "crate" when the feedback shows 'a' was in position 3 for "crate", making the reasoning invalid even if the guess might be acceptable.} \\
\texttt{Wait, let me re{-}evaluate the strategy more carefully.} \\
\texttt{Feedback 1 (slate): a<yellow> (pos 2), t<yellow> (pos 4).} \\
\texttt{Feedback 2 (crate): a<yellow> (pos 3), t<yellow> (pos 4).} \\
\texttt{So 'a' is not in pos 2 or 3. 't' is not in pos 4.} \\
\texttt{The draft explanation says: "Since 'a' was yellow in position 2 (slate) and 3 (crate)..." {-}> This is correct.} \\
\texttt{"...and 't' was yellow in position 4 (slate) and 4 (crate)..." {-}> This is correct.} \\
\texttt{"I will try placing 'a' in position 1 and 't' in position 5..." {-}> This is a valid move.} \\
\texttt{"using common remaining consonants like 'p', 'i', 'n'. 'Pinta' fits these constraints."} \\
\texttt{Is "pinta" a valid word? Yes.} \\
\texttt{Does it avoid s, l, e, c, r? Yes.} \\
\texttt{Does it place 'a' in pos 1? Yes.} \\
\texttt{Does it place 't' in pos 5? Yes.} \\
\texttt{So the strategy seems actually okay?} \\
\texttt{Wait, the prompt says "Do} \\
            }
        }
    }
    & \\ \\

    \theutterance \stepcounter{utterance}  
    & \multicolumn{4}{p{0.6\linewidth}}{
        \cellcolor[rgb]{0.561,0.737,0.561}{
            \makecell[{{p{\linewidth}}}]{
                \texttt{\tiny{[P1$\rangle$GM]}}
                \texttt{explanation: The feedback confirms 'a' and 't' are in the word but in the wrong positions (yellow), while 's', 'l', 'e', 'c', and 'r' are not in the word (red). Since 'a' was yellow in position 2 (slate) and position 3 (crate), it cannot be in those spots. Since 't' was yellow in position 4 in both guesses, it cannot be there. I will try placing 'a' in position 1 and 't' in position 5, using common remaining consonants like 'p', 'i', 'n'. "Pinta" fits these constraints.} \\
\texttt{guess: pinta} \\
            }
        }
    }
    & & \\ \\

    \theutterance \stepcounter{utterance}  
    & & & \multicolumn{2}{p{0.3\linewidth}}{
        \cellcolor[rgb]{0.9,0.9,0.9}{
            \makecell[{{p{\linewidth}}}]{
                \texttt{\tiny{[GM$|$GM]}}
                \texttt{attempts = 3} \\
\texttt{target\_word = faint} \\
\texttt{guess = pinta} \\
\texttt{guess\_feedback = p<red> i<yellow> n<yellow> t<yellow> a<yellow>} \\
            }
        }
    }
    & & \\ \\

    \theutterance \stepcounter{utterance}  
    & & & \multicolumn{4}{p{0.6\linewidth}}{
        \cellcolor[rgb]{0.9,0.9,0.9}{
            \makecell[{{p{\linewidth}}}]{
                \texttt{\tiny{[P1$\langle$GM]}}
                \texttt{guess\_feedback: p<red> i<yellow> n<yellow> t<yellow> a<yellow>} \\
\\ 
\texttt{Provide your response only in this format.} \\
\texttt{explanation: details} \\
\texttt{guess: word} \\
            }
        }
    }
     \\ \\

    \theutterance \stepcounter{utterance}  
    & & \multicolumn{4}{p{0.6\linewidth}}{
        \cellcolor[rgb]{0.95,0.95,0.95}{
            \makecell[{{p{\linewidth}}}]{
                \texttt{\tiny{[P1$\leftrightarrow$P]}}
                \texttt{{[}DRAFT{]}} \\
\texttt{explanation: The feedback shows 'p' is incorrect (red), while 'i', 'n', 't', and 'a' are all in the word but in the wrong positions (yellow). From previous guesses, we know 's', 'l', 'e', 'c', 'r' are also incorrect. The word must consist of the letters \{i, n, t, a\} plus one unknown letter. 'a' was yellow in positions 2, 3, and 5, so it must be in position 1 or 4. 't' was yellow in positions 4 and 5, so it must be in 1, 2, or 3. 'i' was yellow in position 2, so it must be in 1, 3, 4, or 5. 'n' was yellow in position 3, so it must be in 1, 2, 4, or 5. Trying to arrange these: if 'a' is at 1, 't' at 2, 'i' at 3, 'n' at 4, we need a 5th letter. Wait, the word is 5 letters. The letters found are i, n, t, a. We need one more. Let's re{-}evaluate positions. 'a' cannot be 2, 3, 5. 't' cannot be 4, 5. 'i' cannot be 2. 'n' cannot be 3. Let's try "Antic". 'a' at 1 (ok), 'n' at 2 (ok), 't' at 3 (ok), 'i' at 4 (ok), 'c' at 5 (but 'c' was red). So not "Antic". How about "Tania"? 't' at 1, 'a' at 2 (no, 'a' was yellow at 2). "Tina" is 4 letters. "Taint"? 't' at 1, 'a' at 2 (no). "Antit"? No. Let's look at "Tina" + something. Maybe "Tina" is not the root. Let's try "Tania" again? No. What about "Taint"? 't' at 1, 'a' at 2 (bad). "Tina" + 's' (bad). "Tina" + 'l' (bad). "Tina" + 'e' (bad). "Tina" + 'r' (bad). "Tina" + 'c' (bad). "Tina" + 'p' (bad). "Tina" + 'm'? "Tamin"? No. "Tania" is a name. "Taint" has 't' at 1 and} \\
            }
        }
    }
    & \\ \\

    \theutterance \stepcounter{utterance}  
    & & \multicolumn{4}{p{0.6\linewidth}}{
        \cellcolor[rgb]{0.95,0.95,0.95}{
            \makecell[{{p{\linewidth}}}]{
                \texttt{\tiny{[P1$\leftrightarrow$P]}}
                \texttt{{[}SELF{-}FEEDBACK{]}} \\
\texttt{FORMAT: problem} \\
\texttt{RULES: problem} \\
\texttt{STRATEGY: problem} \\
\texttt{REASON: The draft fails to provide a valid guess word, contains internal monologue and incomplete sentences, and violates the strict output format requiring only "explanation:" and "guess:" tags.} \\
            }
        }
    }
    & \\ \\

    \theutterance \stepcounter{utterance}  
    & \multicolumn{4}{p{0.6\linewidth}}{
        \cellcolor[rgb]{0.561,0.737,0.561}{
            \makecell[{{p{\linewidth}}}]{
                \texttt{\tiny{[P1$\rangle$GM]}}
                \texttt{explanation: The feedback confirms 'p' is incorrect, while 'i', 'n', 't', and 'a' are present but misplaced. Previous red letters are 's', 'l', 'e', 'c', 'r'. The word must contain \{i, n, t, a\} plus one unknown letter. Position constraints: 'a' cannot be 2, 3, or 5; 't' cannot be 4 or 5; 'i' cannot be 2; 'n' cannot be 3. Trying "Tania" fails as 'a' is at 2. Trying "Taint" fails as 'a' is at 2. Trying "Antit" is invalid. Let's try "Tina" with a new letter. If 'a' is at 1, 't' at 2, 'i' at 3, 'n' at 4, we need a 5th letter. "Tina" + 'm'? "Tamin"? No. "Tania" is a name. How about "Taint"? No. Let's reconsider "Antic" {-} 'c' is red. "Antia"? 'a' at 1, 'n' at 2, 't' at 3, 'i' at 4, 'a' at 5 (duplicate 'a', but 'a' was yellow in 5, so 'a' cannot be at 5). Wait, 'a' was yellow in 5 in "pinta", so 'a' cannot be at 5. 'a' cannot be 2, 3, 5. So 'a' must be 1 or 4. If 'a' is 1: \_ \_ \_ \_ \_. 't' cannot be 4, 5. 't' can be 2 or 3. 'i' cannot be 2. 'n' cannot be 3. Try "Antia" {-} 'a' at 1, 'n' at 2, 't' at 3, 'i' at 4, 'a' at 5 (invalid, 'a' at 5 is red). Try "Antit" {-} invalid. Try "Tania" {-} 'a' at 2 (invalid). Try "Taint" {-} 'a' at 2 (invalid). Try "Tina" + 'd'? "Tinda"? 't' at 1, 'i' at 2 (invalid, 'i' was yellow at 2). Try "Tina" + 'b'? "Tinda" no. Try "Tina" + 'g'? "Tinga"? 't' at 1, 'i' at 2 (invalid). Try "Tina" + 'h'? "Tinha"? No. Try "T} \\
            }
        }
    }
    & & \\ \\

    \theutterance \stepcounter{utterance}  
    & & & \multicolumn{2}{p{0.3\linewidth}}{
        \cellcolor[rgb]{0.9,0.9,0.9}{
            \makecell[{{p{\linewidth}}}]{
                \texttt{\tiny{[GM$|$GM]}}
                \texttt{Error: The response should contain the 'guess:' keyword exactly once.} \\
            }
        }
    }
    & & \\ \\

    \theutterance \stepcounter{utterance}  
    & & & \multicolumn{2}{p{0.3\linewidth}}{
        \cellcolor[rgb]{0.9,0.9,0.9}{
            \makecell[{{p{\linewidth}}}]{
                \texttt{\tiny{[GM$|$GM]}}
                \texttt{game\_result = ABORT} \\
            }
        }
    }
    & & \\ \\

\end{supertabular}
}

\end{document}
